% !TEX root =  ../Dissertation.tex

\chapter{Conclusion}

This dissertation set out to investigate and improve the performance of Speech Emotion Recognition systems, with a specific focus on enhancing model robustness in noisy environments. Through a systematic series of experiments, we conducted a comparative analysis of multiple feature representations (MFCCs and Mel-spectrograms) and four distinct deep learning architectures: MLP, a shallow CNN, ResNet-18, and the Audio Spectrogram Transformer (AST). The central contribution of this work lies in the extensive evaluation of on-the-fly noise augmentation as a mechanism for improving generalization, using both the RAVDESS and IEMOCAP datasets as testbeds. The main findings and their implications are summarized below.

The experimental results provide conclusive evidence for the architectural superiority of the Audio Spectrogram Transformer. Across all conditions, AST consistently outperformed the other models, achieving the highest accuracy and macro-F1 scores on both clean and noise-augmented test sets. This highlights the effectiveness of its self-attention mechanism in capturing global spectro-temporal dependencies that are crucial for emotion discrimination. Furthermore, on-the-fly noise augmentation proved to be a powerful and effective regularizer. For the AST model, training with a diverse set of natural soundscapes not only fortified its performance against noise but also led to improved results on the original clean data, suggesting that the perturbations encouraged the model to learn more generalizable and salient affective features.


A key insight from this work is the strong interaction between model architecture and the augmentation strategy. The benefits of augmentation were most pronounced for the AST model, which successfully absorbed the acoustic variability to enhance its representations. In contrast, the ResNet-18 architecture, despite its depth, showed fragility and performance degradation when exposed to the same heterogeneous noise, indicating that its convolutional inductive biases may be less suited for disentangling emotional content from complex background sounds. This dissertation also confirmed that performance is highly dependent on dataset characteristics; on the more challenging and spontaneous IEMOCAP dataset, augmentation prompted the AST model to achieve a better precision-recall balance, improving its ability to recognize ambiguous or minority-class emotions.


While this study provides a clear direction for robust SER, several limitations present avenues for future research. The noise sources, though representative, were limited to white noise and a subset of the ESC-50 dataset; future work could incorporate a wider variety of real-world noises, such as babble, music, and channel reverberation. Additionally, while AST proved effective, its performance could be further enhanced by leveraging self-supervised pre-training on large-scale unlabeled datasets, following frameworks like wav2vec 2.0. Another promising direction is the extension of this framework to multimodal SER, integrating textual or visual information, which is particularly relevant for conversational datasets like MELD. Finally, future investigations could explore model interpretability to better understand how attention mechanisms identify and leverage emotionally salient regions within a spectrogram. In summary, the findings of this thesis provide compelling evidence for the adoption of transformer-based architectures combined with robust data augmentation for building the next generation of effective and reliable Speech Emotion Recognition systems.